\subsection*{变量与类型}

Croqtile 引入四类基础类型：标量类型、带形状类型、整数元组类型（\co@ituple@ 类型）和有界类型。它们分别服务于不同的用途。

\begin{itemize}[leftmargin=1.6em]
\item \textbf{标量类型}：为满足程序控制的需要而设计，由\emph{整数类型}和\emph{布尔类型}组成。
\item \textbf{带形状类型}：表示参与计算的数据类型（通常是张量）。除了引用原始数据，\emph{带形状类型}还把数据与一个用多维范围表示的形状关联起来，这种「多维范围」就叫 \co@mdspan@。它对分块等操作很有用。
\item \textbf{整数元组类型}：表示一组整数值，常见用法是给出多维数据的界。
\item \textbf{有界（整数／ITuple）类型}：这是一类特殊类型，用于简化对数据子区域的引用，并与循环构造配合，迭代地处理数据。
\end{itemize}

四类之中，\emph{标量}和 \emph{ituple} 比较容易理解，因为它们对应现有通用编程语言里就有的要素；\emph{带形状类型}和\emph{有界类型}则是这个领域特有的。
